\section{Sorts and Abstract State}

\subsection{Sorts}

The signature contains the following principal sorts:

\begin{align*}
\Bool      &= \{\mathsf{false},\mathsf{true}\},\\
\Result    &= \{\Success,\Failure,\Cancelled\},\\
\Action    &= \{\mathsf{prep},\mathsf{restore},\mathsf{add},\mathsf{flush}\},\\
\Verbosity &= \{\Quiet,\Verbose\}.
\end{align*}

The remaining sorts are characterized rather than exhaustively enumerated:

\begin{description}
  \item[$\DNSName$] A syntactically valid fully qualified DNS name.
  \item[$\Line$] A finite sequence of characters.
  \item[$\Hosts$] A finite sequence of hosts-file lines.
  \item[$\AllowList$] A finite sequence of DNS-name lines.
  \item[$\Backup$] Either $\None$ or $\Some{H}$ for some $H\in\Hosts$.
  \item[$\Metadata$] The mode, owner ID, and group ID of the active hosts file.
  \item[$\Trace$] A finite sequence of diagnostic messages.
  \item[$\Cache$] An abstract state of the DNS cache and
        \code{mDNSResponder}.
\end{description}

\Hosts{} and \AllowList{} are sequences rather than sets because their on-disk
ordering is observable.  A separate set interpretation is used when only
membership matters.

\subsection{State}

The abstract state is the product

\[
  \State = \Hosts \times \Backup \times \AllowList \times
           \Backup \times \Metadata \times \Cache.
\]

A state is written

\[
  S=(H,O,A,B,M,K),
\]

where:

\begin{longtable}{@{}>{$}l<{$}p{0.76\linewidth}@{}}
\toprule
H & active hosts-file contents \\
O & contents of \path{hosts-ORIG}, or $\None$ when absent \\
A & hblock allow-list contents \\
B & contents of \path{hosts.bak}, or $\None$ when absent \\
M & mode and ownership of the active hosts file \\
K & DNS-cache and responder state \\
\bottomrule
\end{longtable}

Field selection is written $S.H$, $S.O$, and so forth.  Real execution also
contains paths, directory entries, timestamps, process identifiers, file
descriptors, and possible faults.  Those details enter the refinement and
failure models rather than every state equation.

\subsection{Execution}

A parsed command produces a result, a new state, and an output trace:

\[
  \Exec : \Command \times \State
       \longrightarrow \Result \times \State \times \Trace.
\]

The trace includes ordinary progress output, verbose metadata, and error
diagnostics.  Separating the trace from the state permits a precise statement
that verbose execution should not change the filesystem result.
